\documentclass[12pt]{report}
\usepackage[utf8x]{inputenc}

%Warning: tipa declares many non-standard macros used by utf8x to
%interpret utf8 characters but extra packages might have to be added
%such as "textgreek" for Greek letters not already in tipa
%or "stmaryrd" for mathematical symbols.
%Utf8 codes missing a LaTeX interpretation can be defined by using
%\DeclareUnicodeCharacter{code}{interpretation}.
%Use coqdoc's option -p to add new packages or declarations.
\usepackage{tipa}

\usepackage[T1]{fontenc}
\usepackage{fullpage}
\usepackage{coqdoc}
\usepackage{amsmath,amssymb}
\usepackage{url}
\begin{document}
%%%%%%%%%%%%%%%%%%%%%%%%%%%%%%%%%%%%%%%%%%%%%%%%%%%%%%%%%%%%%%%%%
%% This file has been automatically generated with the command
%% /nix/store/5nhkcrnn3hzlx1rsgy6zxlg2lrla4f5y-coq-with-packages-9.1.1/bin-orig/coqdoc -latex -toc -interpolate -utf8 -Q . SECF Maps.v Imp.v Equiv.v Hoare.v Hoare2.v Preface.v Noninterference.v StaticIFC.v SpecCT.v Postscript.v Bib.v MapsTest.v ImpTest.v EquivTest.v HoareTest.v Hoare2Test.v PrefaceTest.v NoninterferenceTest.v StaticIFCTest.v SpecCTTest.v PostscriptTest.v BibTest.v 
%%%%%%%%%%%%%%%%%%%%%%%%%%%%%%%%%%%%%%%%%%%%%%%%%%%%%%%%%%%%%%%%%
\coqlibrary{SECF.Postscript}{Library }{SECF.Postscript}

\begin{coqdoccode}
\end{coqdoccode}
\section{Postscript}


\begin{coqdoccode}
\coqdocemptyline
\end{coqdoccode}
\section{Looking Back}



 Here is a quick summary of the topics we covered in this volume: 

\subsection{Noninterference}



\begin{itemize}
\item  definitions for pure functions, state transformers,
  and imperative programs

\item  termination-insensitive noninterference (TINI)

\item  termination-sensitive noninterference (TSNI) 
\end{itemize}


\subsection{Secure multi-execution}



\begin{itemize}
\item  sound and transparent dynamic enforcement of TINI 
\end{itemize}


\subsection{Information-flow control type systems}



\begin{itemize}
\item  type-checkers enforcing TINI and TSNI

\item  for imperative programs with state and outputs 
\end{itemize}


\subsection{Side channels}



\begin{itemize}
\item  control flow security and type system enforcing it

\item  cryptographic constant-time security and type system enforcing it 
\end{itemize}


\subsection{Speculative execution attacks}



\begin{itemize}
\item  speculative constant-time security definition

\item  speculative load hardening (SLH) transformation achieving it 

\end{itemize}
\begin{coqdoccode}
\coqdocemptyline
\end{coqdoccode}
\section{Looking Around}



 The topics above have found practical applications in system security. Below
    we highlight a few recent research projects involving machine-checked proofs: 
\begin{coqdoccode}
\coqdocemptyline
\end{coqdoccode}
\subsection{Proving Noninterference by Parametricity}



 While in \coqdoclibrary{StaticIFC} we showed how to build specialized type systems
    for noninterference, research has shown that in functional programming
    languages with strong type-abstraction mechanisms, information-flow control
    can be implemented as a library. Recent research has shown that in this
    setting simple and elegant noninterference proofs can be built by relying on
    the theory of parametricity, both for libraries doing static enforcement
    \coqdocvar{Algehed} \coqexternalref{and}{http://rocq-prover.org/doc/V9.1.1/stdlib//Corelib.Init.Logic}{\coqdocinductive{and}} \coqdocvar{Bernardy} 2019 (in Bib.v) and for ones doing dynamic enforcement
    \coqdocvar{Algehed} \coqdocvar{et} \coqdocvar{al} 2021 (in Bib.v). These noninterference proofs have been
    machine-checked in Agda. 
\begin{coqdoccode}
\coqdocemptyline
\end{coqdoccode}
\subsection{Formal verification of a constant-time preserving C compiler}



 This work \coqdocvar{Barthe} \coqdocvar{et} \coqdocvar{al} 2020 (in Bib.v) shows that a mildly modified version of
    the CompCert verified C compiler preserves cryptographic constant-time
    security. In particular the authors prove in Rocq that the compiler does
    not introduce secret dependencies into control flow or memory access. Their
    Rocq formalization is aimed at maximizing reuse of the CompCert correctness
    proof, through the use of novel proof techniques for constant-time
    preservation. 
\begin{coqdoccode}
\coqdocemptyline
\end{coqdoccode}
\subsection{Jasmin programming language and compiler}



 Jasmin is a low-level domain-specific language for implementing
    high-assurance and high-speed cryptography. Jasmin programs can be verified
    for correctness, cryptographic security, and side-channel resistance by
    translation to the EasyCrypt proof assistant \coqdocvar{Almeida} \coqdocvar{et} \coqdocvar{al} 2020 (in Bib.v).
    The Jasmin compiler was formally verified in Rocq to be correct
    \coqdocvar{Almeida} \coqdocvar{et} \coqdocvar{al} 2020 (in Bib.v) and to preserve constant-time security
    \coqdocvar{Barthe} \coqdocvar{et} \coqdocvar{al} 2021 (in Bib.v). In more recent work a core compiler inspired by
    Jasmin was proved in Rocq to also preserve speculative constant-time
    \coqdocvar{Arranz}-\coqdocvar{Olmos} \coqdocvar{et} \coqdocvar{al} 2025 (in Bib.v). 
\begin{coqdoccode}
\coqdocemptyline
\end{coqdoccode}
\subsection{Flexible Mechanized Speculative Load Hardening}



 The \coqdoclibrary{SpecCT} chapter and the two projects above are
    aimed at achieving security for cryptographic code. Yet Spectre attacks
    are also a serious threat for non-cryptographic code, since without any
    defenses attackers can construct ``universal read gadgets'' that leak a
    sensitive program's entire memory. SLH is, however, not strong enough for
    protecting code that does not respect the constant-time discipline, leading
    to the introduction of Ultimate SLH \coqdocvar{Zhang} \coqdocvar{et} \coqdocvar{al} 2023 (in Bib.v), which provides
    protection for arbitrary programs, but has too large overhead for general
    use, since it conservatively assumes that all data is secret. More recent
    work introduces Flexible SLH \coqdocvar{Baumann} \coqdocvar{et} \coqdocvar{al} 2025 (in Bib.v), which achieves the
    best of both worlds by generalizing both the selective SLH variant from
    \coqdoclibrary{SpecCT} and Ultimate SLH. Baumann et al prove in Rocq that Flexible
    SLH and Ultimate SLH satisfy a relative security property: any
    transformed program running with speculation must not leak more than what
    the source program leaks sequentially. Their Rocq formalization originated
    as an extension of the simple development from the \coqdoclibrary{SpecCT} chapter. 
\begin{coqdoccode}
\coqdocemptyline
\end{coqdoccode}
\subsection{Strong Timing Isolation of Hardware Enclaves}



 This work \coqdocvar{Lau} \coqdocvar{et} \coqdocvar{al} 2024 (in Bib.v) introduced a RISC-V processor design that is
    formally verified in Rocq to achieve strong timing isolation for enclaves,
    which is formalized in terms of ``air-gaped machines''. 
\begin{coqdoccode}
\coqdocemptyline
\end{coqdoccode}
\section{Looking Forward}



 For readers interesting in research, here are the main conferences
    publishing papers on formal foundations on security:

\begin{itemize}
\item  Computer Security Foundations (CSF)

\item  Principles of Programming Languages (POPL)

\item  International Conference on Functional Programming (ICFP)

\item  Certified Programs and Proofs (CPP)

\item  Interactive Theorem Proving (ITP)

\item  Computer and Communications Security (CCS)

\begin{itemize}
\item  Formal Methods and Programming Languages track

\end{itemize}

\item  IEEE Security and Privacy (SP)

\item  Principles of Secure Compilation Workshop (PriSC) 

\end{itemize}
\begin{coqdoccode}
\coqdocemptyline
\end{coqdoccode}
\end{document}
